\documentclass{amsart}
% For dealing with references we use the comment environment
\usepackage{verbatim}
\newenvironment{reference}{\comment}{\endcomment}
%\newenvironment{reference}{}{}
\newenvironment{slogan}{\comment}{\endcomment}
\newenvironment{history}{\comment}{\endcomment}

% For commutative diagrams we use Xy-pic
\usepackage[all]{xy}

% We use 2cell for 2-commutative diagrams.
\xyoption{2cell}
\UseAllTwocells

% We use multicol for the list of chapters between chapters
\usepackage{multicol}

% This is generally recommended for better output
\usepackage{lmodern}
\usepackage[T1]{fontenc}

% For cross-file-references

% Package for hypertext links:
\usepackage{hyperref}

% For any local file, say "hello.tex" you want to link to please

% Theorem environments.
%
\theoremstyle{plain}
\newtheorem{theorem}[subsection]{Theorem}
\newtheorem{proposition}[subsection]{Proposition}
\newtheorem{lemma}[subsection]{Lemma}

\theoremstyle{definition}
\newtheorem{definition}[subsection]{Definition}
\newtheorem{example}[subsection]{Example}
\newtheorem{exercise}[subsection]{Exercise}
\newtheorem{situation}[subsection]{Situation}

\theoremstyle{remark}
\newtheorem{remark}[subsection]{Remark}
\newtheorem{remarks}[subsection]{Remarks}

\numberwithin{equation}{subsection}

% Macros
%
\def\lim{\mathop{\mathrm{lim}}\nolimits}
\def\colim{\mathop{\mathrm{colim}}\nolimits}
\def\Spec{\mathop{\mathrm{Spec}}}
\def\Hom{\mathop{\mathrm{Hom}}\nolimits}
\def\Ext{\mathop{\mathrm{Ext}}\nolimits}
\def\SheafHom{\mathop{\mathcal{H}\!\mathit{om}}\nolimits}
\def\SheafExt{\mathop{\mathcal{E}\!\mathit{xt}}\nolimits}
\def\Sch{\mathit{Sch}}
\def\Mor{\mathop{\mathrm{Mor}}\nolimits}
\def\Ob{\mathop{\mathrm{Ob}}\nolimits}
\def\Sh{\mathop{\mathit{Sh}}\nolimits}
\def\NL{\mathop{N\!L}\nolimits}
\def\CH{\mathop{\mathrm{CH}}\nolimits}
\def\proetale{{pro\text{-}\acute{e}tale}}
\def\etale{{\acute{e}tale}}
\def\QCoh{\mathit{QCoh}}
\def\Ker{\mathop{\mathrm{Ker}}}
\def\Im{\mathop{\mathrm{Im}}}
\def\Coker{\mathop{\mathrm{Coker}}}
\def\Coim{\mathop{\mathrm{Coim}}}

% Boxtimes
%
\DeclareMathSymbol{\boxtimes}{\mathbin}{AMSa}{"02}

%
% Macros for moduli stacks/spaces
%
\def\QCohstack{\mathcal{QC}\!\mathit{oh}}
\def\Cohstack{\mathcal{C}\!\mathit{oh}}
\def\Spacesstack{\mathcal{S}\!\mathit{paces}}
\def\Quotfunctor{\mathrm{Quot}}
\def\Hilbfunctor{\mathrm{Hilb}}
\def\Curvesstack{\mathcal{C}\!\mathit{urves}}
\def\Polarizedstack{\mathcal{P}\!\mathit{olarized}}
\def\Complexesstack{\mathcal{C}\!\mathit{omplexes}}
% \Pic is the operator that assigns to X its picard group, usage \Pic(X)
% \Picardstack_{X/B} denotes the Picard stack of X over B
% \Picardfunctor_{X/B} denotes the Picard functor of X over B
\def\Pic{\mathop{\mathrm{Pic}}\nolimits}
\def\Picardstack{\mathcal{P}\!\mathit{ic}}
\def\Picardfunctor{\mathrm{Pic}}
\def\Deformationcategory{\mathcal{D}\!\mathit{ef}}

\setlength{\emergencystretch}{3em}
\begin{document}
\title[Stacks Project, Tag 09ZG]{283 --- Global section rigidity for spectral spaces with connected specialization intersections}
\maketitle
\noindent Copyright (C) 2005--2025 Johan de Jong. Stacks Project authors. Source: \href{https://stacks.math.columbia.edu/tag/09ZG}{Tag 09ZG}.

\noindent Editorial difficulty note (not author text). Difficulty H2: Surjectivity of restriction for arbitrary sheaves must be proved from connected intersections with closures of points. The proof reduces via filtered constructible presentations and embeddings into finite closed-supported constant sheaves, while proving that global membership in a subsheaf is detected on the closed subset..

\noindent Editorial provenance note (not author text). History: excerpt from revision \texttt{a0444\allowbreak 6e57e\allowbreak c1fbc\allowbreak 252a8\allowbreak 71afc\allowbreak ec775\allowbreak 2fb28\allowbreak 07b14}, etale-cohomology.tex, lines 15098--15164. Reference presentation was adapted; original-author context and core support blocks were retained or added. The mathematical wording is unchanged. Licensed under GNU FDL 1.2 or later; no invariant sections or cover texts. See ../licenses/COPYING. Context blocks reproduce author definitions and supporting results; they are not additional counted problems.

\begin{definition}
\label{topology-definition-spectral-space}
A topological space $X$ is called {\it spectral} if it is sober,
quasi-compact, the intersection of two quasi-compact opens is
quasi-compact, and the collection of quasi-compact opens forms a
basis for the topology. A continuous map $f : X \to Y$ of spectral
spaces is called {\it spectral} if the inverse image of a quasi-compact
open is quasi-compact.
\end{definition}

\begin{lemma}
\label{lemma-h0-topological}
Let $Z \subset X$ be a closed subset of a topological space $X$.
Assume
\begin{enumerate}
\item $X$ is a spectral space
(Topology, Definition \ref{topology-definition-spectral-space}), and
\item for $x \in X$ the intersection $Z \cap \overline{\{x\}}$
is connected (in particular nonempty).
\end{enumerate}
Then for any sheaf $\mathcal{F}$ on $X$ we have
$\Gamma(X, \mathcal{F}) = \Gamma(Z, \mathcal{F}|_Z)$.
\end{lemma}

\begin{proof}
If $x \leadsto x'$ is a specialization of points, then there is a
canonical map $\mathcal{F}_{x'} \to \mathcal{F}_x$ compatible with
sections over opens and functorial in $\mathcal{F}$. Since every point
of $X$ specializes to a point of $Z$ it follows that
$\Gamma(X, \mathcal{F}) \to \Gamma(Z, \mathcal{F}|_Z)$ is injective.
The difficult part is to show that it is surjective.

\medskip\noindent
Denote $\mathcal{B}$ be the set of all quasi-compact opens of $X$.
Write $\mathcal{F}$ as a filtered colimit $\mathcal{F} = \colim \mathcal{F}_i$
where each $\mathcal{F}_i$ is as in
Modules, Equation (\href{https://stacks.math.columbia.edu/tag/0CAJ}{equation towards constructible sets (Tag 0CAJ)}).
See Modules, Lemma \href{https://stacks.math.columbia.edu/tag/0CAI}{lemma filtered colimit constructibles (Tag 0CAI)}.
Then $\mathcal{F}|_Z = \colim \mathcal{F}_i|_Z$ as restriction to $Z$
is a left adjoint (Categories, Lemma \href{https://stacks.math.columbia.edu/tag/0038}{lemma adjoint exact (Tag 0038)} and
Sheaves, Lemma \href{https://stacks.math.columbia.edu/tag/008K}{lemma f map (Tag 008K)}).
By Sheaves, Lemma \href{https://stacks.math.columbia.edu/tag/009F}{lemma directed colimits sections (Tag 009F)}
the functors $\Gamma(X, -)$ and $\Gamma(Z, -)$ commute with filtered colimits.
Hence we may assume our sheaf $\mathcal{F}$ is as in
Modules, Equation (\href{https://stacks.math.columbia.edu/tag/0CAJ}{equation towards constructible sets (Tag 0CAJ)}).

\medskip\noindent
Suppose that we have an embedding $\mathcal{F} \subset \mathcal{G}$.
Then we have
$$
\Gamma(X, \mathcal{F}) =
\Gamma(Z, \mathcal{F}|_Z) \cap \Gamma(X, \mathcal{G})
$$
where the intersection takes place in $\Gamma(Z, \mathcal{G}|_Z)$.
This follows from the first remark of the proof because we can check
whether a global section of $\mathcal{G}$ is in $\mathcal{F}$ by looking
at the stalks and because every point of $X$ specializes to a point of $Z$.

\medskip\noindent
By Modules, Lemma \href{https://stacks.math.columbia.edu/tag/0CAL}{lemma constructible in constant (Tag 0CAL)}
there is an injection $\mathcal{F} \to \prod (Z_i \to X)_*\underline{S_i}$
where the product is finite, $Z_i \subset X$ is closed, and $S_i$ is finite.
Thus it suffices to prove surjectivity for the sheaves
$(Z_i \to X)_*\underline{S_i}$. Observe that
$$
\Gamma(X, (Z_i \to X)_*\underline{S_i}) = \Gamma(Z_i, \underline{S_i})
\quad\text{and}\quad
\Gamma(X, (Z_i \to X)_*\underline{S_i}|_Z) =
\Gamma(Z \cap Z_i, \underline{S_i})
$$
Moreover, conditions (1) and (2) are inherited by $Z_i$; this is clear
for (2) and follows from
Topology, Lemma \href{https://stacks.math.columbia.edu/tag/0902}{lemma spectral sub (Tag 0902)} for (1). Thus it
suffices to prove the lemma in the case of a (finite) constant sheaf.
This case is a restatement of Lemma \href{https://stacks.math.columbia.edu/tag/0CAM}{lemma connected topological (Tag 0CAM)}
which finishes the proof.
\end{proof}
\end{document}
